\import{algebra/ring/regular-module.tex}
\import{algebra/ring/power.tex}
\import{nat/binomial.tex}
\import{set/logical-graph.tex}

\section{The binomial theorem}

The binomial expansion is a finite sum in the additive group of the ring.
Its summand is displayed directly in the definition, so there is no second
public term for the same expression.

\begin{definition}\label{ring_binomial_expansion}
    $\ringBinomialExpansion{R}{x}{y}{n}
        =\finiteSum{\AdditiveGroup{R}}{\suc{n}}{
            \Apply{\graphOnOp}{
                \suc{n},\carrier[R],
                \fun{k:\Set}{
                    \ringMul[R](
                        \ringNatural{R}{\naturalChoose{n}{k}},
                        \ringMul[R](
                            \ringPower{R}{x}{
                                \naturalDifference{n}{k}},
                            \ringPower{R}{y}{k}))}}}$.
\end{definition}

\begin{proposition}\label{ring_binomial_summand_graph_type}
    Let $R$ be a commutative ring and $x,y\in\carrier[R]$.
    Let $n$ be a natural number.
    Then
        $\Apply{\graphOnOp}{
                \suc{n},\carrier[R],
                \fun{k:\Set}{
                    \ringMul[R](
                        \ringNatural{R}{\naturalChoose{n}{k}},
                        \ringMul[R](
                            \ringPower{R}{x}{
                                \naturalDifference{n}{k}},
                            \ringPower{R}{y}{k}))}}
            \in\funs{\suc{n}}{\carrier[R]}$.
\end{proposition}
\begin{proof}
    $R$ is a ring by \obv{commutative_ring_is_ring}.
    Show that for every $k\in\suc{n}$ we have
        $\ringMul[R](
            \ringNatural{R}{\naturalChoose{n}{k}},
            \ringMul[R](
                \ringPower{R}{x}{\naturalDifference{n}{k}},
                \ringPower{R}{y}{k}))\in\carrier[R]$.
    \begin{subproof}
        Fix $k\in\suc{n}$.
        $k$ is a natural number by
            \auto{omega_successor,naturals_is_transitive_set,
                transitiveset_iff_subseteq,subseteq}.
        $k$ and $n$ are ordinals by
            \obv{natural_number_is_ordinal}.
        $k=n$ or $k\in n$ by \obv{suc_elim}.
        Hence $k\subseteq n$ by
            \obv{subseteq_refl,ordinal_elem_implies_subseteq}.
        $\naturalDifference{n}{k}$ is a natural number by
            \auto{natural_difference_specification}.
        $\naturalChoose{n}{k}$ is a natural number by
            \obv{natural_choose_type}.
        $\ringNatural{R}{\naturalChoose{n}{k}},
            \ringPower{R}{x}{\naturalDifference{n}{k}},
            \ringPower{R}{y}{k}\in\carrier[R]$ by
            \obv{ring_natural_type,ring_power_type}.
        Follows by \obv{ring_mul_type}.
    \end{subproof}
    Follows by \obv{logical_graph_on_function}.
\end{proof}

\begin{lemma}\label{ring_binomial_first_term}
    Let $R$ be a commutative ring and $x,y\in\carrier[R]$.
    Let $n$ be a natural number.
    Then
        $\ringMul[R](
            \ringNatural{R}{\naturalChoose{n}{\emptyset}},
            \ringMul[R](
                \ringPower{R}{x}{\naturalDifference{n}{\emptyset}},
                \ringPower{R}{y}{\emptyset}))
            =\ringPower{R}{x}{n}$.
\end{lemma}
\begin{proof}
    $R$ is a ring by \obv{commutative_ring_is_ring}.
    $\naturalChoose{n}{\emptyset}=\naturalOne$ by
        \obv{natural_choose_first}.
    $\naturalDifference{n}{\emptyset}=n$ by
        \obv{natural_difference_zero}.
    $\ringNatural{R}{\naturalOne}=\ringOne[R]$ by
        \auto{natural_one,ring_natural_one}.
    $\ringPower{R}{y}{\emptyset}=\ringOne[R]$ by
        \obv{ring_power_zero}.
    $\ringPower{R}{x}{n}\in\carrier[R]$ by
        \obv{ring_power_type}.
    Follows by
        \auto{ring_one_type,ring_mul_one_left,ring_mul_one_right}.
\end{proof}

\begin{lemma}\label{ring_binomial_last_term}
    Let $R$ be a commutative ring and $x,y\in\carrier[R]$.
    Let $n$ be a natural number.
    Then
        $\ringMul[R](
            \ringNatural{R}{\naturalChoose{n}{n}},
            \ringMul[R](
                \ringPower{R}{x}{\naturalDifference{n}{n}},
                \ringPower{R}{y}{n}))
            =\ringPower{R}{y}{n}$.
\end{lemma}
\begin{proof}
    $R$ is a ring by \obv{commutative_ring_is_ring}.
    $\naturalChoose{n}{n}=\naturalOne$ by
        \obv{natural_choose_last}.
    $\naturalDifference{n}{n}=\emptyset$ by
        \obv{natural_difference_self}.
    $\ringNatural{R}{\naturalOne}=\ringOne[R]$ by
        \auto{natural_one,ring_natural_one}.
    $\ringPower{R}{x}{\emptyset}=\ringOne[R]$ by
        \obv{ring_power_zero}.
    $\ringPower{R}{y}{n}\in\carrier[R]$ by
        \obv{ring_power_type}.
    Follows by
        \auto{ring_one_type,ring_mul_one_left,ring_mul_one_right}.
\end{proof}

\begin{lemma}\label{ring_binomial_pascal_term}
    Let $R$ be a commutative ring and $x,y\in\carrier[R]$.
    Let $n$ be a natural number and $j\in n$.
    Then
        $\ringMul[R](
            \ringNatural{R}{\naturalChoose{\suc{n}}{\suc{j}}},
            \ringMul[R](
                \ringPower{R}{x}{
                    \naturalDifference{\suc{n}}{\suc{j}}},
                \ringPower{R}{y}{\suc{j}}))
        =\ringAdd[R](
            \ringMul[R](x,
                \ringMul[R](
                    \ringNatural{R}{\naturalChoose{n}{\suc{j}}},
                    \ringMul[R](
                        \ringPower{R}{x}{
                            \naturalDifference{n}{\suc{j}}},
                        \ringPower{R}{y}{\suc{j}}))),
            \ringMul[R](y,
                \ringMul[R](
                    \ringNatural{R}{\naturalChoose{n}{j}},
                    \ringMul[R](
                        \ringPower{R}{x}{
                            \naturalDifference{n}{j}},
                        \ringPower{R}{y}{j}))))$.
\end{lemma}
\begin{proof}
    $R$ is a ring by \obv{commutative_ring_is_ring}.
    $j$ is a natural number by
        \obv{naturals_is_transitive_set,
            transitiveset_iff_subseteq,subseteq}.
    Hence $\suc{j}$ is a natural number by \obv{omega_successor}.
    $j$ and $n$ are ordinals by
        \obv{natural_number_is_ordinal}.
    Hence $j\subseteq n$ and $\suc{j}\subseteq n$ by
        \obv{ordinal_elem_implies_subseteq,ordinal_suc_subseteq}.

    Let $a=\ringNatural{R}{\naturalChoose{n}{j}}$.
    Let $b=\ringNatural{R}{\naturalChoose{n}{\suc{j}}}$.
    Let $u=\ringPower{R}{x}{\naturalDifference{n}{j}}$.
    Let $v=\ringPower{R}{x}{\naturalDifference{n}{\suc{j}}}$.
    Let $s=\ringPower{R}{y}{j}$.
    Let $t=\ringPower{R}{y}{\suc{j}}$.

    $\naturalChoose{n}{j}$ and
        $\naturalChoose{n}{\suc{j}}$ are natural numbers by
        \obv{natural_choose_type}.
    $\naturalDifference{n}{j}$ and
        $\naturalDifference{n}{\suc{j}}$ are natural numbers by
        \auto{natural_difference_specification}.
    Hence $a,b,u,v,s,t\in\carrier[R]$ by
        \auto{ring_natural_type,ring_power_type}.

    $\naturalChoose{\suc{n}}{\suc{j}}
        =\naturalChoose{n}{j}+\naturalChoose{n}{\suc{j}}$ by
        \obv{natural_choose_pascal}.
    Hence
        $\ringNatural{R}{\naturalChoose{\suc{n}}{\suc{j}}}
            =\ringAdd[R](a,b)$ by
        \auto{ring_natural_addition}.

    $\naturalDifference{\suc{n}}{\suc{j}}
        =\naturalDifference{n}{j}$ by
        \obv{natural_difference_successors}.
    $\naturalDifference{\suc{n}}{\suc{j}}
        =\suc{\naturalDifference{n}{\suc{j}}}$ by
        \obv{natural_difference_successor_left}.
    Hence $u=\ringMul[R](v,x)$ by
        \auto{ring_power_successor}.
    $t=\ringMul[R](s,y)$ by
        \auto{ring_power_successor}.

    $\ringMul[R](x,\ringMul[R](b,\ringMul[R](v,t)))
        =\ringMul[R](b,\ringMul[R](u,t))$ by
        \auto{ring_mul_associative,
            commutative_ring_mul_commutative,ring_mul_type}.
    $\ringMul[R](y,\ringMul[R](a,\ringMul[R](u,s)))
        =\ringMul[R](a,\ringMul[R](u,t))$ by
        \auto{ring_mul_associative,
            commutative_ring_mul_commutative,ring_mul_type}.
    $\ringMul[R](\ringAdd[R](a,b),\ringMul[R](u,t))
        =\ringAdd[R](
            \ringMul[R](a,\ringMul[R](u,t)),
            \ringMul[R](b,\ringMul[R](u,t)))$ by
        \obv{ring_distributive_right,ring_mul_type}.
    Follows by
        \auto{ring_add_commutative,ring_mul_type}.
\end{proof}

\begin{proposition}\label{ring_binomial_expansion_type}
    Let $R$ be a commutative ring and $x,y\in\carrier[R]$.
    Let $n$ be a natural number.
    Then $\ringBinomialExpansion{R}{x}{y}{n}\in\carrier[R]$.
\end{proposition}
\begin{proof}
    Let $G=\AdditiveGroup{R}$.
    $R$ is a ring by \obv{commutative_ring_is_ring}.
    $G$ is a commutative group by
        \obv{additive_group_is_exact_commutative_group}.
    Hence $G$ is a monoid by
        \obv{commutative_group_is_group,group_is_monoid}.
    For all $u,v\in\carrier[G]$ we have
        $\mul[G](u,v)=\mul[G](v,u)$ by
        \obv{commutative_group_commutative}.
    $\suc{n}$ is finite by
        \obv{omega_successor,natural_number_is_finite}.
    Let $f=\Apply{\graphOnOp}{
        \suc{n},\carrier[R],
        \fun{k:\Set}{
            \ringMul[R](
                \ringNatural{R}{\naturalChoose{n}{k}},
                \ringMul[R](
                    \ringPower{R}{x}{\naturalDifference{n}{k}},
                    \ringPower{R}{y}{k}))}}$.
    $f\in\funs{\suc{n}}{\carrier[R]}$ by
        \auto{ring_binomial_summand_graph_type}.
    Hence $f\in\funs{\suc{n}}{\carrier[G]}$ by
        \auto{additive_group_fields}.
    $\finiteSum{G}{\suc{n}}{f}\in\carrier[G]$ by
        \auto{finite_sum,monoid_finite_fold_canonical}.
    Hence $\finiteSum{G}{\suc{n}}{f}\in\carrier[R]$ by
        \auto{additive_group_fields}.
    Follows by \auto{ring_binomial_expansion}.
\end{proof}

\begin{proposition}\label{ring_binomial_expansion_zero}
    Let $R$ be a commutative ring and $x,y\in\carrier[R]$.
    Then
        $\ringBinomialExpansion{R}{x}{y}{\emptyset}
            =\ringOne[R]$.
\end{proposition}
\begin{proof}
    Let $G=\AdditiveGroup{R}$.
    Let $f=\Apply{\graphOnOp}{
        \suc{\emptyset},\carrier[R],
        \fun{k:\Set}{
            \ringMul[R](
                \ringNatural{R}{\naturalChoose{\emptyset}{k}},
                \ringMul[R](
                    \ringPower{R}{x}{
                        \naturalDifference{\emptyset}{k}},
                    \ringPower{R}{y}{k}))}}$.
    $R$ is a ring by \obv{commutative_ring_is_ring}.
    $G$ is a commutative group by
        \obv{additive_group_is_exact_commutative_group}.
    Hence $G$ is a monoid by
        \obv{commutative_group_is_group,group_is_monoid}.
    For all $u,v\in\carrier[G]$ we have
        $\mul[G](u,v)=\mul[G](v,u)$ by
        \obv{commutative_group_commutative}.
    $\emptyset$ is a natural number by \obv{omega_empty}.
    Hence $f\in\funs{\suc{\emptyset}}{\carrier[R]}$ by
        \auto{ring_binomial_summand_graph_type}.
    Hence $f\in\funs{\suc{\emptyset}}{\carrier[G]}$ by
        \auto{additive_group_fields}.
    $\suc{\emptyset}=\{\emptyset\}$ by
        \obv{setext,suc_iff,singleton_iff,emptyset}.
    $\naturalChoose{\emptyset}{\emptyset}=\naturalOne$ by
        \obv{natural_choose_first}.
    $\naturalDifference{\emptyset}{\emptyset}=\emptyset$ by
        \obv{natural_difference_zero}.
    $\ringNatural{R}{\naturalOne}=\ringOne[R]$ by
        \auto{natural_one,ring_natural_one}.
    $\ringPower{R}{x}{\emptyset}=\ringOne[R]$ and
        $\ringPower{R}{y}{\emptyset}=\ringOne[R]$ by
        \obv{ring_power_zero}.
    $f(\emptyset)=\ringOne[R]$ by
        \auto{logical_graph_on_apply,ring_one_type,
            ring_mul_one_left,ring_mul_one_right}.
    $\finiteSum{G}{\{\emptyset\}}{f}=f(\emptyset)$ by
        \obv{finite_sum,monoid_finite_fold_singleton,
            singleton_iff}.
    Follows by
        \obv{ring_binomial_expansion}.
\end{proof}

\begin{theorem}\label{ring_binomial_expansion_successor}
    Let $R$ be a commutative ring and $x,y\in\carrier[R]$.
    Let $n$ be a natural number.
    Then
        $\ringBinomialExpansion{R}{x}{y}{\suc{n}}
            =\ringMul[R](
                \ringBinomialExpansion{R}{x}{y}{n},
                \ringAdd[R](x,y))$.
\end{theorem}
\begin{proof}
    Let $G=\AdditiveGroup{R}$.
    Let $M=\LeftRegularModule{R}$.
    Let $A=\suc{n}$.
    Let $B=\suc{A}$.
    $R$ is a ring by \obv{commutative_ring_is_ring}.
    $G$ is a commutative group by
        \obv{additive_group_is_exact_commutative_group}.
    Hence $G$ is a monoid by
        \obv{commutative_group_is_group,group_is_monoid}.
    For all $a,b\in\carrier[G]$ we have
        $\mul[G](a,b)=\mul[G](b,a)$ by
        \obv{commutative_group_commutative}.
    $\LeftRegularModule{R}$ is a left module over $R$ by
        \obv{left_regular_module_is_exact}.
    Hence $M$ is a left module over $R$.
    $\carrier[M]=\carrier[R]$ and
        $\scalarMul[M]=\ringMul[R]$ and
        $\ModuleAdditiveGroup{M}=G$ by
        \auto{left_regular_module_carrier,left_regular_module_fields,
            left_regular_module_additive_group}.
    $\carrier[G]=\carrier[R]$ and
        $\mul[G]=\ringAdd[R]$ and
        $\neutral[G]=\ringZero[R]$ by
        \auto{additive_group_fields}.
    $\suc{n}$ and $\suc{\suc{n}}$ are natural numbers by
        \obv{omega_successor}.
    Hence $A$ and $B$ are natural numbers.
    Hence $A$ and $B$ are finite by
        \obv{natural_number_is_finite}.

    Let $f=\Apply{\graphOnOp}{
        A,\carrier[R],
        \fun{k:\Set}{
            \ringMul[R](
                \ringNatural{R}{\naturalChoose{n}{k}},
                \ringMul[R](
                    \ringPower{R}{x}{\naturalDifference{n}{k}},
                    \ringPower{R}{y}{k}))}}$.
    Let $g=\Apply{\graphOnOp}{
        B,\carrier[R],
        \fun{k:\Set}{
            \ringMul[R](
                \ringNatural{R}{\naturalChoose{\suc{n}}{k}},
                \ringMul[R](
                    \ringPower{R}{x}{
                        \naturalDifference{\suc{n}}{k}},
                    \ringPower{R}{y}{k}))}}$.
    $f\in\funs{A}{\carrier[R]}$ and
        $g\in\funs{B}{\carrier[R]}$ by
        \auto{ring_binomial_summand_graph_type,omega_successor}.
    Hence $f\in\funs{A}{\carrier[M]}$ and
        $f\in\funs{A}{\carrier[G]}$ and
        $g\in\funs{B}{\carrier[G]}$.

    Let $u=\Apply{\graphOnOp}{
        A,\carrier[M],
        \fun{k:\Set}{\scalarMul[M](x,f(k))}}$.
    Let $v=\Apply{\graphOnOp}{
        A,\carrier[M],
        \fun{k:\Set}{\scalarMul[M](y,f(k))}}$.
    $u,v\in\funs{A}{\carrier[M]}$ by
        \auto{logical_graph_on_function,module_scalar_type,
            funs_type_apply}.
    Hence $u,v\in\funs{A}{\carrier[G]}$.

    Let $p=\Apply{\graphOnOp}{
        B,\carrier[G],
        \fun{k:\Set}{
            \Apply{\logicalIfSetOp}{
                (k\in A),\scalarMul[M](x,f(k)),\neutral[G]}}}$.
    Let $q=\Apply{\graphOnOp}{
        B,\carrier[G],
        \fun{k:\Set}{
            \Apply{\logicalIfSetOp}{
                (k=\emptyset),\neutral[G],
                \scalarMul[M](y,f(\unions{k}))}}}$.

    Show that $p\in\funs{B}{\carrier[G]}$.
    \begin{subproof}
        Show that for every $k\in B$ we have
            $\Apply{\logicalIfSetOp}{
                (k\in A),\scalarMul[M](x,f(k)),\neutral[G]}
                \in\carrier[G]$.
        \begin{subproof}
            Fix $k\in B$.
            \begin{byCase}
                \caseOf{$k\in A$.}
                    $f(k)\in\carrier[M]$ by
                        \obv{funs_type_apply}.
                    Hence $\scalarMul[M](x,f(k))\in\carrier[M]$ by
                        \obv{module_scalar_type}.
                    Hence $\scalarMul[M](x,f(k))\in\carrier[G]$.
                    Follows by \obv{logical_if_set_true}.
                \caseOf{$k\notin A$.}
                    $\neutral[G]\in\carrier[G]$ by
                        \obv{monoid_is_unital_magma,unitalmagma_type}.
                    Follows by \obv{logical_if_set_false}.
            \end{byCase}
        \end{subproof}
        Follows by \auto{logical_graph_on_function}.
    \end{subproof}

    Show that $q\in\funs{B}{\carrier[G]}$.
    \begin{subproof}
        Show that for every $k\in B$ we have
            $\Apply{\logicalIfSetOp}{
                (k=\emptyset),\neutral[G],
                \scalarMul[M](y,f(\unions{k}))}
                \in\carrier[G]$.
        \begin{subproof}
            Fix $k\in B$.
            \begin{byCase}
                \caseOf{$k=\emptyset$.}
                    $\neutral[G]\in\carrier[G]$ by
                        \obv{monoid_is_unital_magma,unitalmagma_type}.
                    Follows by \obv{logical_if_set_true}.
                \caseOf{$k\neq\emptyset$.}
                    $k$ is a natural number by
                        \auto{naturals_is_transitive_set,
                            transitiveset_iff_subseteq,subseteq}.
                    Take a natural number $j$ such that $k=\suc{j}$ by
                        \obv{natural_zero_or_successor}.
                    $k=A$ or $k\in A$ by
                        \obv{suc_iff}.
                    Show that $j\in A$.
                    \begin{subproof}
                        \begin{byCase}
                            \caseOf{$k=A$.}
                                Hence $\suc{j}=\suc{n}$.
                                Hence $j=n$ by \obv{suc_injective}.
                                $n\in\suc{n}$ by
                                    \obv{suc_intro_self}.
                            \caseOf{$k\neq A$.}
                                Hence $k\in A$.
                                $A$ is an ordinal by
                                    \obv{natural_number_is_ordinal}.
                                $j\in k$ by \obv{suc_intro_self}.
                                Follows by
                                    \obv{ordinal_transitivity}.
                        \end{byCase}
                    \end{subproof}
                    $j$ is \in-transitive by
                        \obv{natural_number_is_ordinal,
                            ordinal_is_transitiveset}.
                    Hence $\unions{k}=j$ by
                        \auto{transitiveset_iff_unions_suc}.
                    $f(j)\in\carrier[M]$ by
                        \obv{funs_type_apply}.
                    Hence $\scalarMul[M](y,f(j))\in\carrier[M]$ by
                        \obv{module_scalar_type}.
                    Hence $\scalarMul[M](y,f(\unions{k}))
                        \in\carrier[G]$.
                    Follows by \obv{logical_if_set_false}.
            \end{byCase}
        \end{subproof}
        Follows by \auto{logical_graph_on_function}.
    \end{subproof}

    Let $h=\monoidPointwiseMul{G}{B}{p}{q}$.
    $h\in\funs{B}{\carrier[G]}$ by
        \auto{monoid_is_magma,monoid_pointwise_mul_type}.
    Show that $g=h$.
    \begin{subproof}
        Show that for every $k\in B$ we have $g(k)=h(k)$.
        \begin{subproof}
            Fix $k\in B$.
            \begin{byCase}
                \caseOf{$k=\emptyset$.}
                    $\emptyset\in\suc{n}$ by
                        \obv{omega_successor,
                            natural_nonzero_contains_zero,
                            suc_neq_emptyset}.
                    Hence $\emptyset\in A$.
                    $p(\emptyset)=\ringMul[R](x,f(\emptyset))$ and
                        $q(\emptyset)=\ringZero[R]$ by
                        \auto{logical_graph_on_apply,
                            logical_if_set_true}.
                    $f(\emptyset)=\ringPower{R}{x}{n}$ by
                        \auto{logical_graph_on_apply,
                            ring_binomial_first_term}.
                    $g(\emptyset)=\ringPower{R}{x}{\suc{n}}$ by
                        \auto{logical_graph_on_apply,
                            ring_binomial_first_term,omega_successor}.
                    $\ringPower{R}{x}{\suc{n}}
                        =\ringMul[R](\ringPower{R}{x}{n},x)$ by
                        \obv{ring_power_successor}.
                    $\emptyset\in B$.
                    $h(\emptyset)=\ringAdd[R](
                        p(\emptyset),q(\emptyset))$ by
                        \auto{monoid_is_magma,
                            monoid_pointwise_mul_apply}.
                    Follows by
                        \auto{ring_power_type,ring_mul_type,
                            ring_add_zero_right,
                            commutative_ring_mul_commutative}.

                \caseOf{$k\neq\emptyset$.}
                    $k$ is a natural number by
                        \auto{naturals_is_transitive_set,
                            transitiveset_iff_subseteq,subseteq}.
                    Take a natural number $j$ such that $k=\suc{j}$ by
                        \obv{natural_zero_or_successor}.
                    $j$ is \in-transitive by
                        \obv{natural_number_is_ordinal,
                            ordinal_is_transitiveset}.
                    Hence $\unions{k}=j$ by
                        \auto{transitiveset_iff_unions_suc}.
                    $k=A$ or $k\in A$ by
                        \auto{suc_iff}.
                    \begin{byCase}
                        \caseOf{$k\in A$.}
                            $k=n$ or $k\in n$ by
                                \auto{suc_iff}.
                            $n$ is an ordinal by
                                \obv{natural_number_is_ordinal}.
                            $j\in k$ by \obv{suc_intro_self}.
                            Hence $j\in n$ by
                                \auto{ordinal_transitivity}.
                            Hence $j\in A$ by \auto{suc_iff}.
                            $\suc{j}\in A$ by \auto.
                            $p(k)=\ringMul[R](x,f(k))$ and
                                $q(k)=\ringMul[R](y,f(j))$ by
                                \auto{logical_graph_on_apply,
                                    logical_if_set_true,
                                    logical_if_set_false}.
                            $f(\suc{j})=\ringMul[R](
                                \ringNatural{R}{
                                    \naturalChoose{n}{\suc{j}}},
                                \ringMul[R](
                                    \ringPower{R}{x}{
                                        \naturalDifference{n}{\suc{j}}},
                                    \ringPower{R}{y}{\suc{j}}))$ by
                                \auto{logical_graph_on_apply}.
                            Hence $f(k)=\ringMul[R](
                                \ringNatural{R}{
                                    \naturalChoose{n}{\suc{j}}},
                                \ringMul[R](
                                    \ringPower{R}{x}{
                                        \naturalDifference{n}{\suc{j}}},
                                    \ringPower{R}{y}{\suc{j}}))$ by
                                \auto.
                            $f(j)=\ringMul[R](
                                \ringNatural{R}{\naturalChoose{n}{j}},
                                \ringMul[R](
                                    \ringPower{R}{x}{
                                        \naturalDifference{n}{j}},
                                    \ringPower{R}{y}{j}))$ by
                                \auto{logical_graph_on_apply}.
                            $g(\suc{j})=\ringMul[R](
                                \ringNatural{R}{
                                    \naturalChoose{\suc{n}}{\suc{j}}},
                                \ringMul[R](
                                    \ringPower{R}{x}{
                                        \naturalDifference{
                                            \suc{n}}{\suc{j}}},
                                    \ringPower{R}{y}{\suc{j}}))$ by
                                \auto{logical_graph_on_apply}.
                            Hence $g(k)=\ringMul[R](
                                \ringNatural{R}{
                                    \naturalChoose{\suc{n}}{\suc{j}}},
                                \ringMul[R](
                                    \ringPower{R}{x}{
                                        \naturalDifference{
                                            \suc{n}}{\suc{j}}},
                                    \ringPower{R}{y}{\suc{j}}))$ by
                                \auto.
                            Hence $g(k)=\ringAdd[R](
                                \ringMul[R](x,f(k)),
                                \ringMul[R](y,f(j)))$ by
                                \auto{ring_binomial_pascal_term}.
                            $h(k)=\ringAdd[R](p(k),q(k))$ by
                                \auto{monoid_is_magma,
                                    monoid_pointwise_mul_apply}.
                            Follows by \auto.

                        \caseOf{$k\notin A$.}
                            Hence $k=A$ by \auto.
                            Hence $j=n$ by \auto{suc_injective}.
                            $n\in\suc{n}$ by \obv{suc_intro_self}.
                            Hence $n\in A$ by \auto.
                            $p(k)=\ringZero[R]$ and
                                $q(k)=\ringMul[R](y,f(n))$ by
                                \auto{logical_graph_on_apply,
                                    logical_if_set_false}.
                            $f(n)=\ringMul[R](
                                \ringNatural{R}{\naturalChoose{n}{n}},
                                \ringMul[R](
                                    \ringPower{R}{x}{
                                        \naturalDifference{n}{n}},
                                    \ringPower{R}{y}{n}))$ by
                                \auto{logical_graph_on_apply}.
                            Hence $f(n)=\ringPower{R}{y}{n}$ by
                                \auto{ring_binomial_last_term}.
                            $A\in\suc{A}$ by \obv{suc_intro_self}.
                            Hence $\suc{n}\in B$ by \auto.
                            $g(\suc{n})=\ringMul[R](
                                \ringNatural{R}{
                                    \naturalChoose{\suc{n}}{\suc{n}}},
                                \ringMul[R](
                                    \ringPower{R}{x}{
                                        \naturalDifference{
                                            \suc{n}}{\suc{n}}},
                                    \ringPower{R}{y}{\suc{n}}))$ by
                                \auto{logical_graph_on_apply}.
                            Hence $g(\suc{n})
                                =\ringPower{R}{y}{\suc{n}}$ by
                                \auto{ring_binomial_last_term,
                                    omega_successor}.
                            Hence $g(k)=\ringPower{R}{y}{\suc{n}}$ by
                                \auto.
                            $\ringPower{R}{y}{\suc{n}}
                                =\ringMul[R](
                                    \ringPower{R}{y}{n},y)$ by
                                \obv{ring_power_successor}.
                            $h(k)=\ringAdd[R](p(k),q(k))$ by
                                \auto{monoid_is_magma,
                                    monoid_pointwise_mul_apply}.
                            Follows by
                                \auto{ring_power_type,ring_mul_type,
                                    ring_add_zero_left,
                                    commutative_ring_mul_commutative}.
                    \end{byCase}
            \end{byCase}
        \end{subproof}
        $g$ and $h$ are functions and
            $\dom{g}=B=\dom{h}$ by
            \auto{funs_is_function,funs_elim}.
        Follows by \auto{funext_on_domain}.
    \end{subproof}

    Show that $\restrl{p}{A}=u$.
    \begin{subproof}
        $\restrl{p}{A},u\in\funs{A}{\carrier[G]}$ by
            \auto{funs_restrl,subseteq_self_suc_intro}.
        Show that for every $k\in A$ we have
            $\apply{\restrl{p}{A}}{k}=u(k)$.
        \begin{subproof}
            Fix $k\in A$.
            $p$ is a function and $A\subseteq\dom{p}$ by
                \auto{funs_is_function,funs_elim,
                    subseteq_self_suc_intro}.
            Hence $\apply{\restrl{p}{A}}{k}=p(k)$ by
                \obv{restrl_of_function_apply}.
            $A\subseteq\suc{A}$ by \obv{subseteq_self_suc_intro}.
            Hence $A\subseteq B$ by \auto.
            Hence $k\in B$ by \obv{elem_subseteq}.
            $p(k)=\scalarMul[M](x,f(k))$ by
                \auto{logical_graph_on_apply,
                    logical_if_set_true}.
            $u(k)=\scalarMul[M](x,f(k))$ by
                \auto{logical_graph_on_apply}.
        \end{subproof}
        Follows by
            \auto{funs_elim,funext_on_domain}.
    \end{subproof}
    $A\notin A$ by \auto{in_irrefl}.
    $\suc{A}=\cons{A}{A}$ by \obv{suc}.
    Hence $B=\cons{A}{A}$ by \auto.
    $A\in B$ by \auto{cons_iff}.
    $p(A)=\Apply{\logicalIfSetOp}{
        (A\in A),\scalarMul[M](x,f(A)),\neutral[G]}$ by
        \auto{logical_graph_on_apply}.
    Hence $p(A)=\neutral[G]$ by
        \auto{logical_if_set_false}.
    Hence
        $\finiteSum{G}{B}{p}
            =\mul[G](\finiteSum{G}{A}{u},\neutral[G])$ by
        \auto{finite_sum,monoid_finite_fold_fresh_adjoin}.
    $\finiteSum{G}{A}{u}\in\carrier[G]$ by
        \auto{finite_sum,monoid_finite_fold_canonical}.
    Hence
        $\finiteSum{G}{B}{p}=\finiteSum{G}{A}{u}$ by
        \auto{monoid_is_unital_magma,unitalmagma_right}.

    $\scalarMul[M](x,\finiteSum{\ModuleAdditiveGroup{M}}{A}{f})
        =\finiteSum{\ModuleAdditiveGroup{M}}{A}{u}$ by
        \auto{module_scalar_finite_sum}.
    Hence
        $\finiteSum{G}{B}{p}
            =\ringMul[R](x,\finiteSum{G}{A}{f})$ by
        \auto.

    $q(\emptyset)=\neutral[G]$ by
        \auto{logical_graph_on_apply,logical_if_set_true,
            natural_nonzero_contains_zero,suc_neq_emptyset}.
    Show that for every $k\in A$ we have $q(\suc{k})=v(k)$.
    \begin{subproof}
        Fix $k\in A$.
        $k$ is a natural number by
            \auto{naturals_is_transitive_set,
                transitiveset_iff_subseteq,subseteq}.
        Hence $\suc{k}\in B$ by
            \obv{natural_successor_in_successor_bound}.
        $\suc{k}\neq\emptyset$ by \auto{suc_neq_emptyset}.
        $k$ is \in-transitive by
            \obv{natural_number_is_ordinal,
                ordinal_is_transitiveset}.
        Hence $\unions{\suc{k}}=k$ by
            \auto{transitiveset_iff_unions_suc}.
        $q(\suc{k})=\scalarMul[M](y,f(k))$ and
            $v(k)=\scalarMul[M](y,f(k))$ by
            \auto{logical_graph_on_apply,logical_if_set_false}.
    \end{subproof}
    Hence
        $\finiteSum{G}{B}{q}=\finiteSum{G}{A}{v}$ by
        \auto{finite_sum,monoid_finite_fold_successor_shift}.
    $\scalarMul[M](y,\finiteSum{\ModuleAdditiveGroup{M}}{A}{f})
        =\finiteSum{\ModuleAdditiveGroup{M}}{A}{v}$ by
        \auto{module_scalar_finite_sum}.
    Hence
        $\finiteSum{G}{B}{q}
            =\ringMul[R](y,\finiteSum{G}{A}{f})$ by
        \auto.

    $\finiteSum{G}{B}{h}
        =\mul[G](\finiteSum{G}{B}{p},
            \finiteSum{G}{B}{q})$ by
        \auto{finite_sum,monoid_finite_fold_pointwise_mul}.
    Hence
        $\finiteSum{G}{B}{g}
        =\ringAdd[R](
            \ringMul[R](x,\finiteSum{G}{A}{f}),
            \ringMul[R](y,\finiteSum{G}{A}{f}))$ by
        \auto.
    $\finiteSum{G}{A}{f}\in\carrier[R]$ by
        \auto{finite_sum,monoid_finite_fold_canonical}.
    $\ringAdd[R](x,y)\in\carrier[R]$ by \obv{ring_add_type}.
    $\ringMul[R](\ringAdd[R](x,y),\finiteSum{G}{A}{f})
        =\ringAdd[R](
            \ringMul[R](x,\finiteSum{G}{A}{f}),
            \ringMul[R](y,\finiteSum{G}{A}{f}))$ by
        \obv{ring_distributive_right}.
    $\ringMul[R](\ringAdd[R](x,y),\finiteSum{G}{A}{f})
        =\ringMul[R](
            \finiteSum{G}{A}{f},\ringAdd[R](x,y))$ by
        \obv{commutative_ring_mul_commutative}.
    Hence
        $\finiteSum{G}{B}{g}
            =\ringMul[R](
                \finiteSum{G}{A}{f},\ringAdd[R](x,y))$ by
        \auto.
    Follows by \auto{ring_binomial_expansion}.
\end{proof}

\begin{theorem}\label{ring_binomial_theorem}
    Let $R$ be a commutative ring and $x,y\in\carrier[R]$.
    Let $n$ be a natural number.
    Then
        $\ringPower{R}{\ringAdd[R](x,y)}{n}
            =\ringBinomialExpansion{R}{x}{y}{n}$.
\end{theorem}
\begin{proof}
    $R$ is a ring by \obv{commutative_ring_is_ring}.
    $\ringAdd[R](x,y)\in\carrier[R]$ by \obv{ring_add_type}.
    Let $C=\{m\in\naturals\mid
        \ringPower{R}{\ringAdd[R](x,y)}{m}
            =\ringBinomialExpansion{R}{x}{y}{m}\}$.
    $\ringPower{R}{\ringAdd[R](x,y)}{\emptyset}
        =\ringOne[R]$ by \obv{ring_power_zero}.
    $\ringBinomialExpansion{R}{x}{y}{\emptyset}
        =\ringOne[R]$ by
        \obv{ring_binomial_expansion_zero}.
    Hence $\emptyset\in C$ by \obv{omega_empty}.
    Show that for every $m\in C$ we have $\suc{m}\in C$.
    \begin{subproof}
        Fix $m\in C$.
        Hence $m$ is a natural number and
            $\ringPower{R}{\ringAdd[R](x,y)}{m}
                =\ringBinomialExpansion{R}{x}{y}{m}$.
        $\ringPower{R}{\ringAdd[R](x,y)}{\suc{m}}
            =\ringMul[R](
                \ringPower{R}{\ringAdd[R](x,y)}{m},
                \ringAdd[R](x,y))$ by
            \obv{ring_power_successor}.
        $\ringBinomialExpansion{R}{x}{y}{\suc{m}}
            =\ringMul[R](
                \ringBinomialExpansion{R}{x}{y}{m},
                \ringAdd[R](x,y))$ by
            \obv{ring_binomial_expansion_successor}.
        Hence
            $\ringPower{R}{\ringAdd[R](x,y)}{\suc{m}}
                =\ringBinomialExpansion{R}{x}{y}{\suc{m}}$ by
            \auto.
        $\suc{m}$ is a natural number by \obv{omega_successor}.
    \end{subproof}
    Hence $C$ is an inductive set by \obv{inductive_set}.
    Hence every natural number is an element of $C$ by
        \obv{omega_minimal}.
    Thus $n\in C$.
\end{proof}
